\documentclass[10pt, letterpaper]{article}

%------------------------------------------------------------------------------
%  This file contains Zhaoxiang (Aaron) Feng's curriculum vitae rendered using
%  the RenderCV style.
%------------------------------------------------------------------------------

% Packages:
\usepackage[
    ignoreheadfoot, % set margins without considering header and footer
    top=2 cm,      % separation between body and page edge from the top
    bottom=2 cm,   % separation between body and page edge from the bottom
    left=2 cm,     % separation between body and page edge from the left
    right=2 cm,    % separation between body and page edge from the right
    footskip=1.0 cm % separation between body and footer
    % showframe % uncomment for debugging layout
]{geometry}
\usepackage{titlesec}        % customize section titles
\usepackage{tabularx}        % fixed width tables
\usepackage{array}           % required by tabularx
\usepackage[dvipsnames]{xcolor} % color support
\definecolor{primaryColor}{RGB}{0, 79, 144} % RenderCV primary colour
\usepackage{enumitem}        % list customisation
\usepackage{fontawesome5}    % icon fonts
\usepackage{amsmath}         % mathematics support
\usepackage[
    pdftitle={Zhaoxiang (Aaron) Feng's CV},
    pdfauthor={Zhaoxiang Feng},
    pdfcreator={LaTeX with RenderCV},
    colorlinks=true,
    urlcolor=primaryColor
]{hyperref}                  % hyperlinks and PDF metadata
\usepackage[pscoord]{eso-pic}    % floating text on the page
\usepackage{calc}                % length calculations
\usepackage{bookmark}            % PDF bookmarks
\usepackage{lastpage}            % reference the number of pages
\usepackage{changepage}          % adjust width for entries
\usepackage{paracol}             % multi‑column support
\usepackage{ifthen}              % conditional logic
\usepackage{needspace}           % avoid page breaks after titles
\usepackage{iftex}               % engine detection

% Ensure that the generated PDF is machine‑readable and ATS‑parsable.
\ifPDFTeX
    \input{glyphtounicode}
    \pdfgentounicode=1
    \usepackage[utf8]{inputenc}
    \usepackage[T1]{fontenc}
    \usepackage{lmodern}
\fi

%------------------------------------------------------------------------------
%  Layout and formatting settings
%------------------------------------------------------------------------------
\AtBeginEnvironment{adjustwidth}{\partopsep0pt} % remove extra space before adjustwidth
\pagestyle{empty}            % suppress headers and footers
\setcounter{secnumdepth}{0}  % remove section numbering
\setlength{\parindent}{0pt} % no paragraph indentation
\setlength{\topskip}{0pt}   % no top skip
\setlength{\columnsep}{0cm} % spacing between paracol columns

% Customise the section headings to include a horizontal rule beneath.
\titleformat{\section}{\needspace{4\baselineskip}\bfseries\large}{}{0pt}{}[\vspace{1pt}\titlerule]
\titlespacing{\section}{-1pt}{0.3 cm}{0.2 cm}

% Define bullet appearance for the highlights environment.
\renewcommand\labelitemi{$\circ$}
\newenvironment{highlights}{
    \begin{itemize}[
        topsep=0.10 cm,
        parsep=0.10 cm,
        partopsep=0pt,
        itemsep=0pt,
        leftmargin=0.4 cm + 10pt
    ]
}{\end{itemize}}
\newenvironment{highlightsforbulletentries}{
    \begin{itemize}[
        topsep=0.10 cm,
        parsep=0.10 cm,
        partopsep=0pt,
        itemsep=0pt,
        leftmargin=10pt
    ]
}{\end{itemize}}

% Single and two column entry environments.
\newenvironment{onecolentry}{
    \begin{adjustwidth}{0.2 cm + 0.00001 cm}{0.2 cm + 0.00001 cm}
}{\end{adjustwidth}}

\newenvironment{twocolentry}[2][]{
    \onecolentry
    \def\secondColumn{#2}
    \setcolumnwidth{\fill, 4.5 cm}
    \begin{paracol}{2}
}{
    \switchcolumn \raggedleft \secondColumn
    \end{paracol}
    \endonecolentry
}

% Header environment for the CV title and personal details.
\newenvironment{header}{
    \setlength{\topsep}{0pt}\par\kern\topsep\centering\linespread{1.5}
}{
    \par\kern\topsep
}

% Fix for the missing \AND command (renders as a separator pipe)
\newcommand{\AND}{\unskip\enspace\textbar\enspace}

% Save the original \href command and redefine \href to omit arrow icons.
\let\hrefWithoutArrow\href
\renewcommand{\href}[2]{\hrefWithoutArrow{#1}{\ifthenelse{\equal{#2}{}}{ }{#2}}}

\begin{document}

    %--------------------------------------------------------------------------
    %  Header
    %--------------------------------------------------------------------------
    \begin{header}
        \textbf{\fontsize{24 pt}{24 pt}\selectfont Zhaoxiang (Aaron) Feng}

        \vspace{0.3 cm}

        \normalsize
        % The \AND macro inside the header environment draws separators.
        \mbox{{\color{black}\footnotesize\faMapMarker*}\hspace*{0.13cm}San Diego, CA}%
        \kern 0.25 cm% spacing between columns
        \AND%
        \kern 0.25 cm%
        \mbox{\hrefWithoutArrow{mailto:zhf004@ucsd.edu}{\color{black}{\footnotesize\faEnvelope[regular]}\hspace*{0.13cm}zhf004@ucsd.edu}}%
        \kern 0.25 cm%
        \AND%
        \kern 0.25 cm%
        \mbox{\hrefWithoutArrow{tel:+13363377139}{\color{black}{\footnotesize\faPhone*}\hspace*{0.13cm}(336)\,337\,-7139}}%
        \kern 0.25 cm%
        \AND%
        \kern 0.25 cm%
        \mbox{\hrefWithoutArrow{https://github.com/aaronzhfeng}{\color{black}{\footnotesize\faGithub}\hspace*{0.13cm}aaronzhfeng}}%
        \kern 0.25 cm%
        \AND%
        \kern 0.25 cm%        
        \mbox{\hrefWithoutArrow{https://aaronzhfeng.github.io/}{\color{black}{\footnotesize\faGlobe}\hspace*{0.13cm}aaronzhfeng.github.io}}%

    \end{header}

    \vspace{0.3 cm - 0.3 cm}

    %--------------------------------------------------------------------------
    %  Education Section
    %--------------------------------------------------------------------------
    \section{Education}

    % UC San Diego education entry
    \begin{twocolentry}{\textit{Sep. 2022 -- Jun. 2026 (expected)}}
        \textbf{University of California San Diego (UCSD)}

        \textit{B.S. in Data Science \& B.S. in Probability \& Statistics}
    \end{twocolentry}

    \vspace{0.10 cm}
    \begin{onecolentry}
        \begin{highlights}
            \item GPA: 3.83 (Major GPAs: 3.95 in Data Science, 3.93 in Statistics)
            \item Honors: Provost Honors (multiple quarters)
            \item Graduate-level Coursework: Advanced Time Series Analysis, Machine Learning, Applied Statistics, Privacy-sensitive Data Systems
        \end{highlights}
    \end{onecolentry}

    %--------------------------------------------------------------------------
    %  Publication Section
    %--------------------------------------------------------------------------
    \section{Publications}
    
    \begin{onecolentry}
        [1] Matthew Ho, Chen Si, \textbf{Zhaoxiang Feng}, Fangxu Yu, Yichi Yang, Zhijian Liu, Zhiting Hu, Lianhui Qin. ``ArcMemo: Abstract Reasoning Composition with Lifelong LLM Memory.''\\
        \textit{Under review at International Conference on Learning Representations (ICLR)}, 2026. \hfill \textit{[\href{https://arxiv.org/abs/2509.04439}{paper}]} \textit{[\href{https://github.com/matt-seb-ho/arc_memo}{code}]}\\
        \textit{\href{https://arcprize.org/blog/arc-prize-2025-results-analysis}{\textbf{Runner Up}}, ARC Prize 2025 Paper Awards --- Top 8 of 90 paper submissions.}
    \end{onecolentry}

    \vspace{0.20 cm}

    \begin{onecolentry}
        [2] \textbf{Zhaoxiang Feng}, David Scott Lewis, Enrique Zueco. ``LabMemo: Concept-Level Memory for Autonomous Scientific Discovery.''\\
        \textit{Accepted at IEEE IROS 2025 Workshop on Embodied AI and Robotics for Future Scientific Discovery (AIR4S) (non-archival)}, 2025. \hfill \textit{[\href{https://github.com/aaronzhfeng/lab_memo/blob/main/paper/LabMemo_Paper.pdf}{paper}]} \textit{[\href{https://github.com/aaronzhfeng/lab_memo}{code}]}
    \end{onecolentry}

    \vspace{0.20 cm}

    \begin{onecolentry}
        [3] \textbf{Zhaoxiang Feng}, David Scott Lewis. ``SOKRATES: Distilling Symbolic Knowledge into Option-Level Reasoning via Solver-Guided Preference Optimization.''\\
        \textit{Accepted at AAAI 2026 Bridge Program on Logic \& AI: Logical and Symbolic Reasoning in Language Models (LMReasoning)}, 2026. \hfill \textit{[\href{https://github.com/aaronzhfeng/sokrates-oak/blob/main/paper/sokrates.pdf}{paper}]} \textit{[\href{https://github.com/aaronzhfeng/sokrates-oak}{code}]}
    \end{onecolentry}

    %--------------------------------------------------------------------------
    %  Research Experience Section
    %--------------------------------------------------------------------------
    \section{Research Experience}

    % Q-Lab research position
    \begin{twocolentry}{\textit{Jan. 2025 -- Present}}
        \textbf{Research Assistant | Q-Lab, UC San Diego}\\
        \textit{Advisors: Prof. Lianhui Qin and Matthew Ho}
    \end{twocolentry}
    \vspace{0.10 cm}
    \begin{onecolentry}
        \begin{highlights}
            \item Contributed to ArcMemo by leading the design of a program-synthesis-style memory ontology, developing many-to-one puzzle-to-feature mappings and manually curating concept parameterizations to enable concept-level reasoning.
            \item Engineered a reasoning-based retrieval mechanism (System-2 exploration) to resolve embedding failures, achieving a 7.5\% relative gain on ARC-AGI-1 (59.33\% official score).
            \item Built a complete concept dataset generation pipeline transforming hand-written concepts into validated helper puzzles through multi-stage LLM-based generation, code synthesis, and automated testing.
            \item Extended the framework to AIME math problems by designing a metacognitive self-assessment pipeline, improving accuracy by 9.3\% via self-reflective memory usage.
        \end{highlights}
    \end{onecolentry}
    \vspace{0.20 cm}

    % Wang Lab Student Researcher
    \begin{twocolentry}{\textit{Jun. 2025 -- Present}}
        \textbf{Research Assistant | Wang Lab, UC San Diego}\\
        \textit{Advisors: Prof. Wei Wang and Young Su Ko}
    \end{twocolentry}
    \vspace{0.10 cm}
    \begin{onecolentry}
        \begin{highlights}
            \item Developing ProteinAgent (Ongoing): An agentic framework orchestrating domain tools (AlphaFold, MD simulations) using LabMemo-style concepts to guide protocol design for protein stability analysis.
            \item Architected a heterogeneous Mixture-of-Experts system for chemical reaction prediction, integrating four specialized expert models (graph-based, SMILES sequence, 3D geometry-aware, and condition-aware) with learned routing mechanisms on USPTO reaction datasets.
            \item Implemented graph neural network encoders using directed message-passing architectures with shortest-path positional encodings, enabling permutation-invariant molecular representations for stereochemistry-sensitive reaction modeling.
            \item Engineered training pipelines with teacher-forcing, load-balancing losses, and router warmup; designed evaluation frameworks for top-k accuracy and per-expert ablation.
        \end{highlights}
    \end{onecolentry}
    \vspace{0.20 cm}

    % HDSI Capstone Project
    \begin{twocolentry}{\textit{Sep. 2025 -- Present}}
        \textbf{Student Researcher | Halıcıo\u{g}lu Data Science Institute, UC San Diego}\\
        \textit{Advisor: Prof. Hao Zhang}
    \end{twocolentry}
    \vspace{0.10 cm}
    \begin{onecolentry}
        \begin{highlights}
            \item Senior Capstone: "Open LLM Training, Inference, and Infrastructure." Building scalable pipelines for training open-source large language models, benchmarking inference frameworks, and exploring alignment strategies.
        \end{highlights}
    \end{onecolentry}
    \vspace{0.20 cm}

    % USIPC Research Assistant
    \begin{twocolentry}{\textit{Sep. 2023 -- Jun. 2024}}
        \textbf{Research Assistant | U.S. Immigration Policy Center, UC San Diego}\\
        \textit{Advisors: Prof. Tom K. Wong and Dr. Gabriel De Roche}
    \end{twocolentry}
    \vspace{0.10 cm}
    \begin{onecolentry}
        \begin{highlights}
            \item Utilized Python and R to analyze historical naturalization data and forecast long-term immigration trends, informing policy briefs for the 2024 election cycle.
            \item Developed data pipelines and visual dashboards to communicate findings to non-technical stakeholders.
        \end{highlights}
    \end{onecolentry}

    %--------------------------------------------------------------------------
    %  Teaching Experience Section
    %--------------------------------------------------------------------------
    \section{Teaching Experience}

    % Summer Session I 2025 — DSC 40A
    \begin{twocolentry}{\textit{Summer 2025}}
        \textbf{Tutor — DSC\,40A: Theoretical Foundations of Data Science\,I}\\
        University of California San Diego
    \end{twocolentry}

    \vspace{0.10 cm}
    \begin{onecolentry}
        \begin{highlights}
            \item Conducted tutoring sessions covering set theory, probability, and algorithmic thinking for early‑career data science students.
            \item Led review sessions and prepared supplementary materials to aid student learning.
            \item Supervised Oral Exam
        \end{highlights}
    \end{onecolentry}

    \vspace{0.15 cm}

    % Grader positions (condensed)
    \begin{twocolentry}{\textit{Fall 2024 -- Fall 2025}}
        \textbf{Grader — Mathematics Department}\\
        University of California San Diego
    \end{twocolentry}

    \vspace{0.10 cm}
    \begin{onecolentry}
        \begin{highlights}
            \item Graded for MATH\,180A/C (Probability \& Stochastic Processes), MATH\,181A (Mathematical Statistics), and MATH\,185 (Computational Statistics) a total of 7 times.
        \end{highlights}
    \end{onecolentry}

    %--------------------------------------------------------------------------
    %  Technical Skills Section
    %--------------------------------------------------------------------------
    \section{Technical Skills}

    \begin{onecolentry}
        \textbf{Languages:} Python, Java, R, SQL, JavaScript, HTML\\[3pt]
        \textbf{ML Frameworks:} PyTorch, PyTorch Geometric, scikit-learn, Transformers, NVIDIA PhysicsNeMo, neural operators (Fourier Neural Operator)\\[3pt]
        \textbf{LLM/NLP:} OpenAI APIs, prompt engineering, LLM fine-tuning (PPO, DPO), RAG systems, memory‑augmented LLMs, reasoning‑based retrieval, vision‑language models\\[3pt]
        \textbf{Data \& Infrastructure:} Pandas, Dask, Spark, AWS, Docker, distributed training, vector databases (embedding retrieval)\\[3pt]
        \textbf{Dev Tools:} Git/GitHub, Linux, Hydra, SSH\\[3pt]
        \textbf{Scientific ML \& Agents:} physics‑based ML (PDE/inverse problems), continual learning (Elastic Weight Consolidation), concept‑level memory design, autonomous scientific agents, safety‑aware Planner/Selector/Verifier architectures
    \end{onecolentry}

\end{document}